Phase Windings as Directional Propagation of Topology  
The Dual Residual Structures Generated by a Single Seed

Topology, in the present setting, is not an abstract invariant imposed after the fact upon a finished geometric figure. It is the oriented propagation of lines—their successive directions, their cumulative turning, and the residual that remains when that turning fails to close upon an exact rational multiple of the circle. Phase windings are precisely this directional propagation. They are the continuous record of how orientation itself advances, and the geometry that finally appears is the global organisation of those oriented lines. Two complementary operations, both generated by one residual seed, convert that oriented propagation into the two primary residual structures of the hierarchy: the residual bulge and the residual hexagon.

1. Orientation as Primary Topological Datum

A phase winding is an oriented path on the circle (or on the unit sphere) whose instantaneous direction is the argument of a complex (or hypercomplex) phase. The maximal winding \(\theta_{\max}=6.5\pi\) is such a path: its total turning is three and a quarter full rotations, and its asymptotic terminal ray carries the definite orientation \(\arctan(2\pi)\). The direction of this ray is not an accessory label; it is the topological residue of the entire winding. Every subsequent geometric construction inherits its orientation from this directional datum.

2. The Residual Seed as Oriented Defect

When the continuous angular bridge \(\psi\) is iterated seven times, the accumulated orientation falls short of (or overshoots) the exact six-fold lattice by the residual

\[
\delta=7\psi-\pi/3\approx0.0051676144\ \mathrm{rad}.
\]

This \(\delta\) is an oriented defect. It measures the failure of the directional propagation to close upon a pure rational symmetry. Because the defect is itself an oriented quantity, every structure that descends from it remains marked by a preferred sense of rotation—the same counterclockwise-positive sense fixed by the original imaginary unit. The residual seed is therefore not a scalar error; it is a topological orientation that refuses to vanish.

3. First Operation: Winding Method and the Residual Bulge

The winding method follows the oriented lines of the maximal winding until they confront the altitude of the residual triangle. The linear gap

\[
1.490969476641-1.412965136507=0.078004340134
\]

is the directional distance between the terminal ray and the height of the near-equilateral cell generated by the same residual seed. From that cell the leftover area of the unit disk is extracted, partitioned, and ratioed to produce the dimensionless intensity \(\rho\). The intensity is then returned to the boundary as a normal displacement of each oriented edge:

\[
\gamma_\rho(t)=\gamma_0(t)+c\cdot\rho\cdot4t(1-t)\,n.
\]

The resulting bulge is the continuous geometric expression of the original oriented defect. In topological terms, the winding method converts directional propagation into a curvature of the very lines that carry the orientation. The residual bulge is oriented topology made visible as edge deformation.

4. Second Operation: Radial-Symmetry Folding and the Residual Hexagon

The same oriented lines, once the residual seed has acted six times, are compelled to close. Closure under an oriented residual cannot produce a perfect regular hexagon; it produces a practical hexagon whose sides still carry the directional memory of \(\delta\). The diagonals of this hexagon recover two interlocking residual triangles; their mutual intersections generate a second, inner hexagon. Both the outer practical hexagon and the inner intersection hexagon are therefore organisations of oriented lines—global arrangements forced by the requirement that the directional propagation return to itself after six residual steps. Radial-symmetry folding converts the same oriented defect into a discrete closed form. The residual hexagon is oriented topology made visible as combinatorial closure.

5. Dual Unity under a Single Seed

Because both operations begin from the identical residual seed, the bulge and the hexagon are not independent inventions. They are dual readings of one oriented propagation. The winding method follows the lines far enough to extract an intensity and return it as curvature; the folding method binds the lines into a closed hexagonal skeleton and an inner hexagonal measure. The curvature deforms the sides of the very triangles whose intersections produce the inner hexagon; the inner hexagon, in turn, sits inside the curved figure as its discrete topological core. Orientation propagates through both structures without interruption.

6. Hierarchical Consequence

Every higher densification merely multiplies the number of oriented lines while preserving their residual intensity and their residual sense of rotation. The hierarchy is therefore the progressive enrichment of a single directional propagation: more lines, the same orientation, the same residual defect, the same dual expression as curvature and as hexagonal closure. Topology, in this precise sense, is the history of how orientation advances, fails to close exactly, and organises itself into the residual bulge and the residual hexagon—the two primary residual structures generated by one and the same oriented seed.
